\chapter{Related Work}
\label{ch:related-work}

This thesis builds on three lines of work: the systematization of polynomial hash implementation techniques, program generation for fast prime-field polynomial hashing, and newer multivariate polynomial hash designs motivated by higher-security applications.

\section{Efficient Prime-Field Polynomial Hashing}
\label{sec:sok}

Degabriele, Gilcher, Govinden, and Paterson give a systematic overview of design and implementation choices for polynomial hash functions over prime fields~\cite{degabriele2024sok}. They separate high-level design choices, such as the polynomial construction and finite field, from implementation choices, such as limb representation, multiplication strategy, and carry propagation. This separation is central to our compiler: the DSL describes the polynomial structure, while compiler settings select the field, target platform, vectorization strategy, multiplication algorithm, and carry schedule.

Their work also emphasizes the importance of delayed carry propagation and delayed limb realignment. Unsaturated limb representations leave enough space in each machine word to accumulate temporary overflow bits. This can reduce the number of sequential carry chains, which are costly because they limit instruction-level parallelism. Our compiler exposes this idea through explicit carry instructions in the IR and through a configurable realignment mode.

The paper also proposes new prime-field designs that revisit Poly1305's design choices on modern hardware. These designs are important context for this thesis because they demonstrate that changing the prime field and the implementation strategy can improve performance, security, or both.

\section{SPHGen}
\label{sec:sphgen}

SPHGen is a program generator for fast polynomial hash functions over Crandall prime fields~\cite{pegolotti2026sphgen}. It takes field and target parameters as input and emits optimized implementations for vector ISAs. Its generated programs contain specialized load, add, multiply, carry, and main-loop code for the selected field representation and target architecture.

Our implementation uses a similar code-generation setting, but exposes it through a small DSL rather than through fixed polynomial templates. We use Crandall prime fields, unsaturated limb representations, vectorized arithmetic routines, and configurable multiplication algorithms, while letting the source program describe the algebraic structure to be lowered.

The difference is the level of abstraction exposed to the programmer. SPHGen targets a specific class of simple polynomial hash functions. The project in this thesis instead starts from a DSL that can express several patterns.

\section{New Multivariate Polynomial Designs}
\label{sec:newdesigns}

The work on new multivariate-polynomial universal hash functions proposes two-level polynomial designs that combine lower-level and higher-level hash functions~\cite{gilcher2026newdesigns}. At the lower level, the paper studies MMH, SQH, NMH, and HKM constructions. The precise definitions used as reference examples are collected in Appendix~\ref{app:polynomial-definitions}. These functions are particularly relevant for this thesis as they serve both as propelling examples and as test subjects for our compiler.

MMH is naturally expressed as a sum of products. SQH replaces products by squarings, incentivizing a dedicated square operation. NMH groups message and key elements pairwise and is also naturally expressed as a sum. HKM is different: it is a sequential recurrence. The definition in the paper appears to contain a typo in the recursive call; throughout this thesis we use the intended recurrence. For generated code, a direct recursive implementation is unattractive, because it introduces function-call overhead and prevents the compiler from controlling carry placement. We therefore express the same recurrence iteratively using a left fold.

\section{Prime-Field Hashing Beyond Poly1305}
\label{sec:prime-field-beyond-poly1305}

The workshop paper on universal hash designs discusses why polynomial hash functions should not be tied to the exact Poly1305 parameter set~\cite{degabriele2024accordion}. Its broader message is that prime-field polynomial hashing remains attractive beyond \(2^{130}-5\), especially when the desired security level, message-length regime, or target platform differs from the setting for which Poly1305 was originally designed.

Larger fields can give stronger security; smaller or better-aligned fields can improve throughput; and designs with different polynomial structures can trade key size, multiplication count, and parallelism differently. This reinforces the need for tooling that can vary the field and the implementation parameters without rewriting the hash function by hand.
